
\section{Variants} \label{sec:variants}

We now discuss the consequences of imposing the following 
restrictions on symmetric \cm{s}:
% 
\begin{itemize}
\item the set of $\xztinstr$ transitions is also \emph{symmetric}, 
that is invariant under permutations of counters; 
\item all $\xupdinstr$ transitions $(q,\updinstr{w},q')$ are 
\emph{single-counter}, that is $w$  is nonzero on at most one coordinate.
\end{itemize}
% 
\para{Symmetric zero-tests}

For simplicity, let's focus on 2-\cm{s} below.
Symmetric $\xztinstr$ transitions, effectively testing
that an \emph{unspecified} counter is zero,
written $\ztinstr{\ast}$,
can surprisingly simulate unrestricted ones.
To this aim, a configuration $q(v_1,v_2)$ of a symmetric 2-\cm would 
be simulated by $q(2v_1,2v_2+1)$, thus doubling the counter values and 
distinguishing them by parity.
Consequently, a transition $(q,\updinstr{w_1,w_2},q')$ is simulated
a transition $(q,\updinstr{2w_1,2w_2},q')$.
Therefore the $\xupdinstr$ transitions, even symmetric ones, preserve
parity of both counters.
Due to parity distinction, a transition
$(q,\ztinstr{\ast},q')$ checking if an unspecified counter is 0,
in fact checks if counter 1 is 0.
On the other hand,
simulation of transition $(q,\ztinstr 2, q')$
is done in three steps,
using two auxiliary intermediate states $p_t,p'_t$:
\begin{align} \label{eq:sim3}
(q,\updinstr{1,-1},p_t) \qquad
(p_t, \ztinstr{\ast}, p'_t) \qquad
(p'_t, \updinstr{-1,1},q').
\end{align}
Observe that the first and the last transition swap parity of
counters, and likewise do their swapped versions.
Therefore a successful $\ztinstr{\ast}$ is only possible on
counter 2, which rules out possibility of use of swapped versions
of the first and last transition.

\smallskip

The above considerations generalise 
to $G$-\cm{s}, for 
an arbitrary number $d>2$ of counters, and arbitrary permutation group
$G\leq\symg d$.
To this aim we distinguish between the counters by remainders modulo 
$d$, simulate a configuration $q(a_1, \ldots, a_d)$ by 
\[q(d a_1, d a_2 + 1, \ldots, d a_d + d-1),\]
and simulate a transition $(q,\updinstr{\vr v}, q')$ by
$(q,\updinstr{d\vr v}, q')$.
Furthermore, to correctly simulate a zero-test of counter $j+1>1$, 
say a transition $(q,\ztinstr {j+1},q')$,
we replace the vectors $(1,-1)$ and $(-1,1)$ in the first
and last transition in \eqref{eq:sim3} by
\[
(j,\ldots,j,-j, j, \ldots, j)
\qquad\text{ and } \qquad
(-j,\ldots,-j,j, -j, \ldots, -j),
\]
respectively, where the unique value of negative (left) or positive
(right) sign appears on the $(j+1)$th coordinate.
% \mei{I don't think it should be $d$, I expected $j-1$ and $-(j-1)$
% respectively, to match the special case worked out for 2-\cm{s}.}

\para{Symmetric zero-tests and single-counter updates}
We do not know how much weaker is the fragment with only 
single-counter $\xupdinstr$ transitions (but unrestricted
$\xztinstr$ transitions).
However, imposing both restrictions results in a degeneracy,
namely the model is simulated by (and equivalent to) 1-\vass,
at least as long as zero source and target vectors are considered.
Indeed, 
whenever a symmetric \cm has a run 
$q(0,\ldots, 0)\step^* q'(0,\ldots, 0)$, it also has a
\emph{simple} run, in which all counter updates are performed on
counter 1, say, while all zero-tests are (trivially) performed on counter 2,
say.
The zero-tests are thus void, and simple runs are essentially
runs of a 1-\vass.
